\section{Coefficient scales}
\label{app:coordinates}

The parameter cover is expressed in boundary coefficients, so we need
both a local coordinate change from \(\mu_1,\mu_2,\mu_3\) to
\(q_0,q_1,q_2\) and a comparison of their weighted scales. We compute
the coordinate change first and then prove the scale bounds directly.

On \(z=0\), let \(v_{\rm bd}\) be the small root of
\(\mu_3+v_{\rm bd}+v_{\rm bd}^2=0\), and put
\(t=1+2v_{\rm bd}\). Writing \(v=v_{\rm bd}+u\) gives
\[
 \eta=tu+u^2,\qquad
 u=\eta/t-\eta^2/t^3+O(\eta^3).
\]
The boundary identity
\eqref{eq:boundary-drift-conjugacy-new} reads
\[
 q(\eta)=(t+2u)\{P(v_{\rm bd})+P'(v_{\rm bd})u-3v_{\rm bd}u^2-u^3\}.
\]
It follows that
\begin{equation}
\label{eq:direct-coefficient-bridge}
 q_0=tP(v_{\rm bd}),\qquad q_1=P'(v_{\rm bd})+2P(v_{\rm bd})/t,\qquad
 q_2=-3v_{\rm bd}/t+q_1/t^2-4q_0/t^4.
\end{equation}
Differentiating \eqref{eq:direct-coefficient-bridge} at the organizing
point gives the Jacobian and its inverse:
\[
 D_{\mu}(q_0,q_1,q_2)\big|_0
 =\begin{pmatrix}1&0&0\\[1mm]2&1&0\\[1mm]-2&1&3\end{pmatrix},
 \qquad
 D_q(\mu_1,\mu_2,\mu_3)\big|_0
 =\begin{pmatrix}1&0&0\\[1mm]-2&1&0\\[1mm]4/3&-1/3&1/3\end{pmatrix}.
\]
Thus the coefficient change is locally invertible.

Set
\[
 H=\max\{|q_0|^{1/3},\sqrt{|q_1|},|q_2|\},\qquad
 H_\mu=\max\{|\mu_1|^{1/3},\sqrt{|\mu_2|},|\mu_3|\}.
\]
These scales are comparable on a fixed neighborhood.
Indeed \eqref{eq:direct-coefficient-bridge} gives \(|v_{\rm bd}|\le CH\)
when \(1/2<t<3/2\), and the inverse identities
\[
 \mu_2=q_1-2q_0/t^2+3v_{\rm bd}^2,\qquad
 \mu_1=q_0/t-\mu_2v_{\rm bd}+v_{\rm bd}^3,\qquad
 \mu_3=-v_{\rm bd}-v_{\rm bd}^2
\]
give \(H_\mu\le CH\). Conversely \(v_{\rm bd}=O(H_\mu)\),
\(P(v_{\rm bd})=O(H_\mu^3)\), \(P'(v_{\rm bd})=O(H_\mu^2)\), and
\eqref{eq:direct-coefficient-bridge} imply \(H\le CH_\mu\).

\begin{samepage}
For \(H>0\), put \(\bar q_j=q_j/H^{3-j}\). Uniformly on this
compact normalized coefficient set,
\[
 \frac {v_{\rm bd}}H=-\frac{\bar q_2}{3}+O(H),\qquad
 \frac{\mu_3}{H}=\frac{\bar q_2}{3}+O(H),\qquad
 \frac{\mu_2}{H^2}=\bar q_1+\frac{\bar q_2^2}{3}+O(H),
\]
\[
 \frac{\mu_1}{H^3}
   =\bar q_0+\frac{\bar q_1\bar q_2}{3}+\frac{2\bar q_2^3}{27}+O(H),\qquad
 H^{-3}q(HX)=\bar q_0+\bar q_1X+\bar q_2X^2-X^3+O_{C^k}(H).
\]
\end{samepage}
The last estimate, for every fixed \(k\) on every fixed \(X\)-interval,
follows from analyticity and the fact that the boundary drift depends only on
\((\mu_1,\mu_2,\mu_3)\); \(A,C\) remain independent parameters
in the positive-height coefficient \(B\).
